% !TeX program = xelatex
\documentclass[UTF8,fontset=fandol]{ctexart}
\usepackage[a4paper,margin=22mm,top=20mm,bottom=20mm]{geometry}
\usepackage{tabularx,booktabs,hyperref,fancyhdr,xcolor}
\providecommand{\speakername}{负责HyMEM部分的同学}
\providecommand{\nextspeakername}{下一位}
\hypersetup{pdftitle={HyMEM 录音讲稿：统一68页演示版},pdfauthor={\speakername},colorlinks=true,urlcolor=blue}
\pagestyle{fancy}\fancyhf{}
\fancyhead[L]{HyMEM 录音讲稿}\fancyhead[R]{\speakername}
\fancyfoot[C]{\thepage}\setlength{\headheight}{16pt}
\setlength{\parindent}{2em}\setlength{\parskip}{7pt}
\linespread{1.18}
\begin{document}\zihao{-4}\raggedbottom
\begin{center}{\Large\bfseries HyMEM 录音讲稿}\par
统一68页演示：第37--43页\quad 计划8分钟\par
\speakername\quad 2026年10月27日\end{center}
\noindent 本页为录音提示，不念进音频。后续每页对应一张幻灯片，只读正文。
\subsection*{计划翻页安排}
\noindent{\small\begin{tabularx}{\textwidth}{@{}rrrX@{}}\toprule
幻灯片 & 开始 & 结束 & 内容\\\midrule
37 & 00:00 & 00:20 & HyMEM：电脑操作智能体的混合自演化记忆 \\
38 & 00:20 & 01:20 & 网页操作中的经验表示与更新 \\
39 & 01:20 & 03:05 & 混合记忆：策略与轨迹表示 \\
40 & 03:05 & 04:50 & 长期经验与工作记忆的更新 \\
41 & 04:50 & 05:45 & 实验设置与编码器训练 \\
42 & 05:45 & 07:05 & 混合记忆的任务成功率 \\
43 & 07:05 & 08:00 & 方法总结与课设启发 \\
\bottomrule\end{tabularx}}\par
\subsection*{录音时注意}
\begin{itemize}
\item 开头自我介绍；结尾请下一位同学接着讲。全程只需录音，不必录屏。
\item 本表是计划。录完后，从最终音频开头计时，记录每张幻灯片第一句话的实际起点。
\item 页码、时间与标题不用念。数字按中文口语读，适当在流程图和结果表前后停顿。
\item 若剪辑或重录，需按最终音频重填时间点。还没有录音，不能把计划表当作实测记录。
\end{itemize}
\subsection*{术语提示}
HyMEM、CLIP可连贯读出；ADD、MERGE、REPLACE随后直接解释为“新增、合并、替换”。模型名称可读作“千问二点五的七B视觉语言模型”。
\clearpage
\noindent{\small 幻灯片第37页\hfill 计划 00:00--00:20}
\section*{HyMEM：电脑操作智能体的混合自演化记忆}
大家好，我是\speakername{}。接下来我介绍 HyMEM，也就是面向电脑操作智能体的混合自演化记忆。这篇论文主要想回答两个问题：过去的操作经验怎么保存，以及任务进行到一半，经验该怎样更新。
\clearpage
\noindent{\small 幻灯片第38页\hfill 计划 00:20--01:20}
\section*{网页操作中的经验表示与更新}
先看它要解决的任务。假如我们让一个智能体帮忙买一件商品，它需要看懂网页截图，搜索、筛选，再进入结账。过去成功做过的事情当然有参考价值，但保存这些经验并不简单。

一部分经验可以写成策略，比如先按价格筛选，再检查配送条件；另一部分却藏在界面和动作里，例如按钮在哪里、点开之后出现什么页面。只用几句话总结，容易把这些细节丢掉。

而且，经验是否相关还会随着任务变化。搜索时需要的是筛选方法，到了结账就要关注地址和支付；不过，用户的预算仍然需要保留。

所以，HyMEM把问题拆成两层：用混合表示决定存什么，再用记忆更新，让智能体在不同阶段拿到合适的经验。
\clearpage
\noindent{\small 幻灯片第39页\hfill 计划 01:20--03:05}
\section*{混合记忆：策略与轨迹表示}
这一页左边是使用记忆的流程，右边是一条经验包含的三个层次。我们先看右边。

第一层是高层策略，回答的是这类任务通常怎么做。第二层是属性，比如搜索、筛选、价格，帮助把相关经验联系起来。第三层是连续的轨迹表示，保留截图和动作中不容易完整写成文字的细节。公式里的三个量，就对应这三层内容。

回到左边。系统先把任务描述和当前截图编码，用相似度找到五条最相关的经验，作为种子。接着沿着图中的属性关联，再选出五条邻居，最后使用十条经验。

为什么不直接取十条最相似的？因为看起来相似的经验，不一定覆盖完成任务需要的所有步骤。图中的邻居有机会补充关联知识。当然，邻居也可能带来噪声，所以五加五是论文采用的设置，是否适合别的任务还需要比较。

取回经验后，策略和属性会被整理成文本指导，连续轨迹嵌入则与模型输入一起使用。这里有个容易混淆的地方：CLIP生成的向量是用来找到经验的；送给智能体使用的连续嵌入，则负责保留经验细节。它们的作用不同。
\clearpage
\noindent{\small 幻灯片第40页\hfill 计划 03:05--04:50}
\section*{长期经验与工作记忆的更新}
记忆取回来以后，也不能一直原样用下去。HyMEM把更新分成两种：左边是跨任务维护长期经验，右边是调整当前任务的工作记忆。

先看长期更新。新轨迹会先找到相关旧经验，再由视觉语言模型判断，它到底提供了什么新信息。出现新策略或新属性，就执行 ADD，也就是新增；如果策略相同，但补充了另一种界面或操作路径，就执行 MERGE，把互补证据整合进去。

如果新轨迹在相同条件下更好，比如步骤更少，或者执行更稳，就用 REPLACE 更新旧证据。因此，相似只是把候选找出来，是否保留、合并或替换，还需要进一步判断。

右边的局部刷新则发生在任务执行过程中。每次动作之后，模型比较前后观测，判断任务有没有进入新的阶段。比如从搜索进入结账，就保留预算和目标，放下过时的搜索指导，再按照新状态检索，并更新文本指导和连续嵌入。

可以把两种更新理解为：长期经验库记住哪些做法值得复用，工作记忆决定眼下该参考什么。它不是每一步都把全部记忆推倒重来，而是在识别出阶段变化后，保留有用部分，再补充新内容。
\clearpage
\noindent{\small 幻灯片第41页\hfill 计划 04:50--05:45}
\section*{实验设置与编码器训练}
实验是在三个多模态网页操作基准上进行的，名称列在第一行。前两个主要根据截图和操作轨迹判断任务是否完成，MMInA则评价最终回答。表里的Overall，是对可用领域的成功率取平均。

记忆来自两千八百八十三条成功轨迹，主图有一千八百五十八个节点。执行时，还需要界面标记和定位能力配合。

另外，论文会训练记忆编码器中的Q-Former和LoRA，可训练参数大约占百分之一点二，通常训练一个epoch。这与刚才的记忆内容更新是两个阶段。理解结果时，记忆数据、训练和界面定位这些条件，都需要一起考虑。
\clearpage
\noindent{\small 幻灯片第42页\hfill 计划 05:45--07:05}
\section*{混合记忆的任务成功率}
结果表先看第一行，也就是千问二点五的七B视觉语言模型。无记忆时成功率是百分之十二点五，只用离散记忆是百分之十七，只用连续记忆是百分之二十一点七，而HyMEM达到百分之三十五。

从十二点五到三十五，提高的是二十二点五个百分点。这个对照说明，在论文的设置下，把策略指导和轨迹细节结合起来，比只保留其中一部分更有效。

再看另外两行，混合记忆也获得了更高的总体结果，但收益大小并不一样。尤其是千问三这一行，离散记忆只有百分之二十一点三，反而略低于无记忆的百分之二十一点九。

这提醒我们，增加记忆不一定就能提高表现，还要看记忆内容和模型是否配合。表里展示的是论文作者的网页操作实验；对我们的课设，更值得借鉴的是机制，而不是直接比较不同任务的分数。
\clearpage
\noindent{\small 幻灯片第43页\hfill 计划 07:05--08:00}
\section*{方法总结与课设启发}
最后总结一下。HyMEM把策略和视觉细节结合起来，再通过长期经验更新和任务内刷新，让记忆跟着任务状态变化。它的收益也伴随着代价：需要训练编码器，维护图关联，还要在执行时做判断，错误标签或更新判断也可能影响后续操作。

对我们做的文本事件记忆检索，我更关注右边这几个问题：语义种子和结构邻居怎样分工；在固定检索预算下，扩展邻居究竟补充了多少关键证据，又引入了多少噪声和调用成本。

这些问题可以进一步通过对照实验来回答。我们的文本证据检索，与论文直接使用多模态轨迹表示的条件不同，需要分别验证。

我的部分到此结束，下面有请\nextspeakername{}同学继续介绍MetaMem。
\end{document}
